\documentclass[11pt]{article}
\usepackage{amsmath, amssymb}
\usepackage{geometry}
\usepackage{array}
\usepackage{booktabs}
\geometry{margin=0.7in}
\setlength{\parindent}{0pt}
\setlength{\parskip}{0.35em}
\renewcommand{\arraystretch}{1.15}

\newcommand{\cheatsheettitle}[2]{%
\begin{center}
{\LARGE #1}\\[0.4em]
{\large #2}
\end{center}
}


\begin{document}
\cheatsheettitle{Algebra II Functions and Transformations Cheatsheet}{Module 2}

\section*{Parent Functions}
\begin{center}
\begin{tabular}{lll}
\toprule
Family & Parent & Key Feature \\
\midrule
linear & $f(x)=x$ & constant rate \\
quadratic & $f(x)=x^2$ & vertex \\
cubic & $f(x)=x^3$ & origin symmetry \\
absolute value & $f(x)=|x|$ & corner vertex \\
square root & $f(x)=\sqrt x$ & starts at endpoint \\
reciprocal & $f(x)=\frac1x$ & asymptotes \\
exponential & $f(x)=b^x$ & horizontal asymptote \\
logarithmic & $f(x)=\log_b x$ & vertical asymptote \\
\bottomrule
\end{tabular}
\end{center}

\section*{Transformation Form}
\[
g(x)=a\,f(b(x-h))+k
\]
\[
h=\text{horizontal shift},\qquad k=\text{vertical shift}
\]
\[
|a|=\text{vertical stretch/compression},\qquad a<0\Rightarrow \text{reflect over }x\text{-axis}
\]
\[
\text{horizontal scale factor}=\frac1{|b|},\qquad b<0\Rightarrow \text{reflect over }y\text{-axis}
\]
\[
|b|>1\Rightarrow \text{horizontal compression},\qquad 0<|b|<1\Rightarrow \text{horizontal stretch}
\]

\section*{Domain and Range}
\[
\text{domain}=\text{allowed inputs}
\]
\[
\text{range}=\text{possible outputs}
\]
Check denominators, even roots, logarithms, graph endpoints, and asymptotes.

\section*{Composition}
\[
(f\circ g)(x)=f(g(x))
\]
The domain of $f\circ g$ must satisfy:
\[
x\in \text{domain of }g,\qquad g(x)\in \text{domain of }f
\]

\section*{Inverses}
\[
f(f^{-1}(x))=x,\qquad f^{-1}(f(x))=x
\]
To find an inverse:
\begin{enumerate}
  \item Write $y=f(x)$.
  \item Switch $x$ and $y$.
  \item Solve for $y$.
  \item State domain and range restrictions.
\end{enumerate}
Horizontal line test: a function has an inverse function if every horizontal line hits the graph at most once.

\section*{Piecewise Functions}
\[
f(x)=
\begin{cases}
2x+1, & x<0\\
x^2, & x\ge0
\end{cases}
\]
Use the rule whose condition contains the input.

\section*{Quick Checks}
\begin{itemize}
  \item Transform inside the input affects $x$-values.
  \item Transform outside the function affects $y$-values.
  \item Inverse graphs reflect across $y=x$.
  \item Piecewise boundaries control which rule applies.
\end{itemize}

\end{document}
